\subsection{Constructing the geometric prior family}

Let $Y-X-Z$ and let the encoder posterior be
$q_\phi(z|x)=\mathcal N(\mu_\phi(x),\operatorname{diag}\sigma_\phi^2(x))$.
CEB~\citep{fischer2020conditional} replaces the VIB marginal reference by a
label-conditional reference $r(z|y)$ and uses the certificate
\begin{equation}
\E\KL(q_\phi(z|X)\|r(z|Y))
=I(X;Z|Y)+\E_Y\KL(q_\phi(z|Y)\|r(z|Y))\ge I(X;Z|Y).
\label{eq:gap}
\end{equation}
The certificate is a sum of class-wise matching terms. At $r=q(z|y)$ it reduces to
$I(X;Z|Y)$, independently of the pairwise arrangement of the class conditionals
(Proposition~\ref{prop:geom-blind}); the $T_\alpha$ counterexample is deferred to
Appendix~\ref{app:diagnosis}.

\begin{proposition}[The certificate records nothing about inter-label geometry]
\label{prop:geom-blind}
The CEB compression term contains no term coupling two distinct class priors; at its
class-wise matched optimum it equals
\begin{equation}
\mathrm{ReX}=I(X;Z|Y),
\label{eq:reix-residual}
\end{equation}
regardless of their inter-class means or directions. The certificate therefore
measures the amount of residual compression but none of the arrangement of the
class conditionals, so the compression criterion alone cannot prefer any inter-class
geometry.
\end{proposition}

\paragraph{Qualified geometric prior.}
GPB restricts the reference to stateless conditional Gaussians
\begin{equation}
p_\theta(z|y)=\mathcal N(m_\theta(y),\Sigma_\theta(y)),\qquad
\Sigma_\theta(y)\preceq\tau_{\max}^2I,
\label{eq:pgib-covariance}
\end{equation}
whose means obey a declared two-sided separation for every pair of distinct
discrete labels $y\ne y'$,
\begin{equation}
\Delta\le\|m_\theta(y)-m_\theta(y')\|\le\bar\Delta,
\label{eq:pgib-separation-condition}
\end{equation}
or a two-sided label metric condition for $P_Y$-almost every label pair
$(y,y')$, where $d_Y$ is a metric on the label space $\mathcal Y$:
\begin{equation}
\rho_{\min}d_Y(y,y')\le\|m_\theta(y)-m_\theta(y')\|
\le\rho_{\max}d_Y(y,y'),
\label{eq:pgib-lipschitz}
\end{equation}
with constants fixed before training. The objective is
\begin{equation}
\mathcal L_{\rm GPB}=\E[\CE(c_\psi(z),Y)+\beta\KL(q_\phi(z|X)\|p_\theta(z|Y))],
\qquad z\sim q_\phi(z|X).
\label{eq:pgib-objective}
\end{equation}

\paragraph{Instances.}\label{sec:opb}
For classification, OPB evaluates a prior encoder on all class labels and uses its
thin polar factor $Q_\theta$. If $U_\theta$ is the resulting all-class matrix, then
\begin{equation}
Q_\theta=U_\theta(U_\theta^\top U_\theta)^{-1/2},
\qquad Q_\theta^\top Q_\theta=I_K,
\label{eq:opb-qr}
\end{equation}
and we assume, monitoring rather than imposing it, that the monitored matrix has
full column rank (Assumption~\ref{ass:fullrank} in the appendix). Writing
$q_{\theta,k}$ for the $k$-th column of the frame $Q_\theta$, the anchors are
\begin{equation}
p_\theta^{\rm OPB}(z|Y=k)=\mathcal N(aq_{\theta,k},\tau^2I),\qquad
Q_\theta^\top Q_\theta=I_K,\qquad\|aq_i-aq_j\|=\sqrt2a.
\label{eq:opb-prior}
\end{equation}
The construction also gives the fixed geometry
\begin{equation}
(aq_i)^\top(aq_j)=a^2\mathbf 1\{i=j\},
\qquad \|aq_i-aq_j\|=\sqrt2a\quad(i\ne j).
\label{eq:opb-geometry}
\end{equation}
The anchor-energy scores are
\begin{equation}
s_k(z)=-\frac{\|z-aq_{\theta,k}\|^2}{2\tau^2}+\log\pi_k,
\label{eq:opb-energy}
\end{equation}
the Gaussian Bayes scores; with uniform priors this induces an orthogonal classifier
whose scale is fixed by the prior (Appendix~\ref{app:construction}). For OPB, the
general objective is equivalently
\begin{equation}
\mathcal L_{\rm OPB}=\E\left[\CE(c_\psi(Z),Y)+\beta\KL\bigl(q_\phi(z|X)\|p_\theta^{\rm OPB}(z|Y)\bigr)\right],
\qquad Z\sim q_\phi(\cdot|X),
\label{eq:opb-objective}
\end{equation}
so the same stateless anchors enter both the KL and the tied readout. Detailed polar
construction, initialization, covariance relaxation, and the relation to
prototype/ETF methods are given in Appendices~\ref{app:construction} and~\ref{sec:related}.

For regression, EPB uses a one-dimensional continuous label, standardized once on
the training set as $\tilde y$ with the statistics frozen. Let $W\in\R^{d\times1}$ be
a trainable direction with per-step normalization $u=W/\|W\|$; the isometric prior is
\begin{equation}
p_\theta^{\rm EPB}(z|y)=\mathcal N(\rho\tilde y\,u,\tau^2I),\qquad
\|\mu_p(y)-\mu_p(y')\|=\rho|\tilde y-\tilde y'|,
\label{eq:epb-prior}
\end{equation}
so latent distances equal $\rho$ times standardized label distances exactly. Order
and magnitude of label differences are encoded with equality, which qualifies the
family (Eq.~\eqref{eq:pgib-lipschitz} with $\rho_{\min}=\rho_{\max}=\rho$). The
prediction head is the tied projection $\hat{\tilde y}=u^\top z/\rho$, the regression
analogue of the anchor-energy classifier: no free direction or scale can bypass the
prior. The main variant uses the trainable normalized axis with no bias; the
fixed-axis and unconstrained-scale ablations are in
Appendix~\ref{app:construction}.

\begin{proposition}[Stateless GPB priors preserve the CEB certificate]
\label{prop:opb-certificate}
For every conditional prior $p_\theta(z|y)$ independent of $x$,
\begin{equation}
\E\KL\bigl(q_\phi(z|X)\|p_\theta(z|Y)\bigr)
=I(X;Z|Y)+\E_Y\KL\bigl(q_\phi(z|Y)\|p_\theta(z|Y)\bigr)
\ge I(X;Z|Y).
\label{eq:opb-certificate}
\end{equation}
Thus restricting the prior family changes its geometry without changing the CEB
certificate identity or its lower-bound reading, so the geometry constraint is a
free move for compression; batch-dependent references would not satisfy this
statement.
\end{proposition}

\subsection{Alignment and separation guarantees}
\label{sec:guarantees}
Define
\begin{equation}
\mathcal R(w)=\E[\ell(c_\psi(Z),Y)],\qquad
\mathcal K(w)=\E\KL(q_\phi(z|X)\|p_\theta(z|Y)),
\label{eq:risk-rate-definitions}
\end{equation}
and
\begin{equation}
\mathcal D(w)=\frac{\E\|\mu_\phi(X)-m_\theta(Y)\|^2}{2\tau_{\max}^2}.
\label{eq:expected-anchor-mismatch}
\end{equation}
\begin{theorem}[Comparator-dependent alignment under approximate realizability]
\label{prop:pgib-alignment}
If $\widehat w$ is $\varepsilon$-optimal for $\mathcal R+\beta\mathcal K$, then for
every feasible comparator $w_0$,
\begin{equation}
\mathcal D(\widehat w)\le\mathcal K(\widehat w)\le
\mathcal K(w_0)+\frac{\mathcal R(w_0)+\varepsilon}{\beta}.
\label{eq:pgib-comparator-bound}
\end{equation}
Thus a comparator with total KL $\delta_{\rm app}$ and risk $C_{\rm app}$ gives
\begin{equation}
\mathcal D(\widehat w)\le\delta_{\rm app}
 +\frac{C_{\rm app}+\varepsilon}{\beta}.
\label{eq:pgib-approximate-alignment}
\end{equation}
The exact $O(1/\beta)$ rate is only the zero-floor, exactly realizable special case;
the full ambiguity--approximation decomposition is in Appendix~\ref{app:deferred}.
The bound thus certifies alignment to whatever comparator is prespecified: its KL is
the irreducible floor, and the risk term controls how fast the fitted solution
approaches that floor as $\beta$ grows.
\end{theorem}

For a target $d_\star$, the bound is informative only if
$\mathcal K(w_0)<d_\star$ and
\begin{equation}
\beta\ge\frac{\mathcal R(w_0)+\varepsilon}{d_\star-\mathcal K(w_0)}.
\label{eq:comparator-informative-condition}
\end{equation}
A prespecified comparator's KL and risk can be estimated on an independent held-out
set using
\begin{equation}
\widehat{\mathcal K}_S(w_0)=\frac1{|S|}\sum_{(x_i,y_i)\in S}
\KL(q_{\phi_0}(z|x_i)\|p_{\theta_0}(z|y_i)),\quad
\widehat{\mathcal R}_S(w_0)=\frac1{|S|}\sum_{(x_i,y_i)\in S}
\E_{q_{\phi_0}}\ell(c_{\psi_0}(Z),y_i),
\label{eq:heldout-comparator-estimates}
\end{equation}
which estimates a comparator-specific floor, not an intrinsic dataset floor.

\begin{theorem}[Mismatch-conditioned posterior separation]
\label{thm:pgib-separation}
Let
\begin{equation}
D_k=\frac{\E[\|\mu_\phi(X)-m_\theta(k)\|^2\mid Y=k]}{2\tau_{\max}^2}.
\label{eq:pgib-class-rate}
\end{equation}
For discrete labels,
\begin{equation}
\|m_i-m_j\|\ge\Delta-\sqrt{2\tau_{\max}^2D_i}-\sqrt{2\tau_{\max}^2D_j}.
\label{eq:pgib-separation}
\end{equation}
For continuous labels define
\begin{equation}
D(y)=\frac{\E[\|\mu_\phi(X)-m_\theta(y)\|^2\mid Y=y]}{2\tau_{\max}^2},
\label{eq:pgib-class-rate-continuous}
\end{equation}
and the analogous lower bound is
\begin{equation}
\|m(y)-m(y')\|\ge\rho_{\min}d_Y(y,y')
-\sqrt{2\tau_{\max}^2D(y)}-\sqrt{2\tau_{\max}^2D(y')},
\label{eq:pgib-separation-continuous}
\end{equation}
for almost every pair. The discrete bound is positive only when
\begin{equation}
\Gamma_{ij}:=\frac{\sqrt{2\tau_{\max}^2D_i}+\sqrt{2\tau_{\max}^2D_j}}
{\|m_\theta(i)-m_\theta(j)\|}<1;
\label{eq:separation-tightness-ratio}
\end{equation}
for OPB with equal mismatch $D$ this is $D<a^2/(4\tau_{\max}^2)$. The scale is a
validation-selected design parameter and does not itself guarantee a positive bound.
The declared margin thus survives up to the two measured correction terms, so the
guarantee is informative exactly when the measured mismatch stays below the margin;
small measured mismatch transfers the declared geometry to the posterior almost
intact.
\end{theorem}
